\documentclass{article}

\usepackage[margin=1in]{geometry}
\linespread{1.2}
\usepackage{amsmath,amsthm,amssymb,amsfonts}

 
\title{Real Analysis I --- Homework 2}

% Replace with your name and ID below
\author{FirstName LastName \\ ID: 123-45-6789}

\date{}


\begin{document}

\maketitle

Consider $\mathbb{R}^{n}$ and $\mathcal{M} = \mathcal{M}(\mathbb{R}^{n})$. If not specified, $n$ is arbitrary; otherwise $n=1$.
\begin{enumerate}

\item Suppose $E \notin \mathcal{M}$ and $\mu(A) = 0$. Is $E \setminus A \in \mathcal{M}$? Prove your claim.

%uncomment below and write your proof after it.
%\begin{proof}
%Write your proof here.
%\end{proof}


%\newpage % uncomment this line to start on a new page

\item Show that $E \in \mathcal{M}$ iff for any $A \subset E$ and $B \subset E^{c}$, there is $\mu^{*} (A \cup B) = \mu^{*} (A) + \mu^{*} (B)$.

%uncomment below and write your proof after it.
%\begin{proof}
%Write your proof here.
%\end{proof}


%\newpage % uncomment this line to start on a new page

\item Show that $E \in \mathcal{M}$ iff there exists $G_{\delta}$-set $H$ such that $E \subset H$ and $\mu^{*} (H \setminus E) = 0$.

%uncomment below and write your proof after it.
%\begin{proof}
%Write your proof here.
%\end{proof}


%\newpage % uncomment this line to start on a new page

\item Suppose $A \in \mathcal{M}$, and $B$ is an arbitrary set (not necessarily in $\mathcal{M}$). Show that $\mu^{*} (A \cap B) + \mu^{*} (A \cup B) = \mu(A) + \mu^{*} (B)$.

%uncomment below and write your proof after it.
%\begin{proof}
%Write your proof here.
%\end{proof}


%\newpage % uncomment this line to start on a new page

\item Suppose $\mu^{*} (A \triangle B) = \mu^{*} (B \triangle C) = 0$. Show that $\mu^{*} (A \triangle C) = 0$.

%uncomment below and write your proof after it.
%\begin{proof}
%Write your proof here.
%\end{proof}


%\newpage % uncomment this line to start on a new page

\item Suppose $A$ and $B$ are bounded. Show that $| \mu^{*}(A) - \mu^{*}(B) | \le \mu^{*} (A \triangle B)$.

%uncomment below and write your proof after it.
%\begin{proof}
%Write your proof here.
%\end{proof}


%\newpage % uncomment this line to start on a new page


\item Suppose $A \subset [0,1]$ is measurable and $ \mu(A) = 1$. Show that for any measurable $B \subset [0,1]$, there is $\mu (A \cap B) = \mu (B)$.

%uncomment below and write your proof after it.
%\begin{proof}
%Write your proof here.
%\end{proof}


%\newpage % uncomment this line to start on a new page


\item Does there exist a nonempty open set $G$ such that $\mu(G)=0$? Prove your claim.

%uncomment below and write your proof after it.
%\begin{proof}
%Write your proof here.
%\end{proof}


%\newpage % uncomment this line to start on a new page


\item Suppose $G$ is open and $E$ has measure zero. Prove that $\overline{G} = \overline{G \setminus E}$.

%uncomment below and write your proof after it.
%\begin{proof}
%Write your proof here.
%\end{proof}


%\newpage % uncomment this line to start on a new page


\item Does there exist a closed proper subset $F \subsetneq [a,b]$ such that $\mu(F) = b-a$? Prove your claim.

%uncomment below and write your proof after it.
%\begin{proof}
%Write your proof here.
%\end{proof}



\end{enumerate}


\end{document}